\begin{remark}[三类投影在凸资源几何中的统一]\label{rem:pom-convex-projection-unification}
上述凸几何框架表明三类投影并非彼此孤立。规范投影 $P_Z$ 给出对象层的稳定像，原子投影 $P_{\mathrm{prim}}$ 给出轨道原子，序投影 $P_{\le}$ 给出可比较的平均量与阈值门，而多势压力 $P(\theta)$ 及其对偶率函数 $I_{\theta_0}$ 把这三者压进一个连续、可优化且可对偶化的几何壳——即可达均值凸体 $\mathcal{R}_\varphi$ 与压力面 $\theta\mapsto P(\theta)$。因此本体演化 $U$ 不仅生成对象层上的稳定代表与 primitive 原子，还同时生成对象层上可见的\emph{凸资源几何}：不同轨道、不同平衡态与不同权重偏置所能呈现的平均读出组合，恰好由 $\mathcal{R}_\varphi$ 的凸约束与 $P(\theta)$ 的 Legendre--Fenchel 对偶完全记录。
\end{remark}
